\documentclass[11pt]{article}
\usepackage[T1]{fontenc}
\usepackage{amsmath,amssymb,amsthm,mathtools}
\usepackage[margin=1in]{geometry}
\usepackage{xcolor}
\usepackage{booktabs,longtable,array}
\usepackage[hidelinks]{hyperref}

\newcommand{\PROVED}{\textcolor{green!45!black}{\textbf{PROVED}}}
\newcommand{\OPEN}{\textcolor{red!70!black}{\textbf{OPEN}}}
\newcommand{\AUDIT}{\textcolor{orange!80!black}{\textbf{AUDIT}}}
\newcommand{\CONDITIONAL}{\textcolor{blue!65!black}{\textbf{CONDITIONAL}}}
\newcommand{\Ran}{\operatorname{Ran}}
\newcommand{\sgn}{\operatorname{sgn}}
\newcommand{\old}{\mathrm{old}}
\newcommand{\new}{\mathrm{new}}

\title{Row (d), G40, and the Two Distinct \(d\)-Constructions\\
       Corrected Full-Closure Audit (v131)}
\author{}
\date{August 14, 2026}

\begin{document}
\maketitle

\begin{center}
\fbox{\parbox{0.92\textwidth}{
\textbf{Controlling verdict.}
The old local-analysis row (d) and the new Gamma--Fourier--Tate row (d) are
different constructions and are audited separately below.  The files presently
prove substantial exact identities, reductions, and local positivity statements,
but they do not prove G40, condition \(G\), or row (d) for every window \(T\).
Those three statements therefore remain \OPEN.  No numerical endpoint, Krein
factorization, definition of a sign operator, or conditional cascade is counted
as the missing positivity theorem.
}}
\end{center}

\section{Scope: the two objects called row (d)}

\subsection{The old row (d)}

The old object is the local analytic construction contained in
\texttt{main.tex} and \texttt{row-d-local-analysis.tex}.  Its relevant proved
content is:
\begin{enumerate}
\item the explicit-formula operator and its primitive compression are defined;
\item full Hilbert-space positivity is rigorously certified for
      \(0<T\leq \log 2\);
\item at the next finite threshold, in particular near
      \(T=\tfrac12\log 5\), finite-dimensional and complementary estimates are
      obtained separately;
\item the audit identifies the missing joint Schur coupling and stops there.
\end{enumerate}
Thus the old row (d) supplies a rigorous base interval, finite-threshold tests,
and the correct Schur/capacity obstruction.  It does not supply global positivity
for all \(T\).

\subsection{The new row (d)}

The new object is the external sequence v89--v130.  It develops a
Gamma--Fourier--Tate boundary/current model, a moving Tate cell, cross-only
transport, a Gamma source state, rational leakage, and several candidate closure
mechanisms.  Its authoritative accumulated operator ledger is v94; the later
files refine particular routes without replacing the v94 terminal target.

\medskip
\noindent\AUDIT\quad A theorem proved for the old local model may be used as an
input to the new construction only through an explicitly proved comparison map.
Conversely, a formal or conditional identity in the new construction does not
retroactively complete the old local proof.

\section{The frozen target: no further movement of the goal}

For a threshold cell, the new construction has the data
\[
 A_0=\mathcal G_0^*\mathcal G_0\geq0,
 \qquad q=M_0h',
 \qquad \beta=S_M'.
\]
The threshold block and its shorted defect are
\begin{equation}
 \mathcal K_T=
 \begin{pmatrix}A_0&q\\q^*&\beta\end{pmatrix},
 \qquad
 \Delta_T=\beta-q^*A_0^\dagger q .
 \tag{F1}
\end{equation}
The range condition is
\begin{equation}
 q\in\Ran A_0^{1/2}.
 \tag{F2}
\end{equation}
The direct G40 target in the energy notation of v94 is
\begin{equation}
 S_E-Z_E^*Z_E-b_E^*A_0^\dagger b_E\geq0,
 \qquad b_E\in\Ran A_0^{1/2}.
 \tag{F3}
\end{equation}
Equivalently, the desired complete-source factorization is
\begin{equation}
 \mathcal K_T=\mathcal T_T^*\mathcal T_T
 \quad\text{in a Hilbert space, for every threshold and every }T.
 \tag{F4}
\end{equation}

The v94 equivalence theorem correctly establishes
\begin{equation}
 \mathcal K_T\geq0
 \Longleftrightarrow
 \bigl[(\mathrm{F2})\text{ and }\Delta_T\geq0\bigr]
 \Longleftrightarrow
 J_T:=\sgn(\Delta_T)=I
 \Longleftrightarrow G40
 \Longleftrightarrow G
 \tag{F5}
\end{equation}
in the stated threshold framework.

\medskip
\noindent\PROVED\quad The definitions, Moore--Penrose Schur theorem, and the
equivalences in (F5).

\noindent\OPEN\quad The positive side of any one of these equivalent statements
for every \(T\).  Defining \(\Delta_T\) and \(J_T\), or constructing a Krein
factorization with signature \(J_T\), does not prove \(\Delta_T\geq0\) or
\(J_T=I\).

\section{What rows (a), (b), and (c) actually contribute}

\begin{longtable}{>{\raggedright\arraybackslash}p{0.10\textwidth}
                  >{\raggedright\arraybackslash}p{0.35\textwidth}
                  >{\raggedright\arraybackslash}p{0.43\textwidth}}
\toprule
Row & Proved contribution & Boundary of the proved statement \\
\midrule
\endfirsthead
\toprule
Row & Proved contribution & Boundary of the proved statement \\
\midrule
\endhead
(a) & \PROVED\ in the intrinsic periodic--Deligne--nuclear category; it constructs
the required nuclear correspondence framework. & It does not give a finite-rank
realization of the full kernel; the manuscript proves an obstruction to such a
realization. \\
\addlinespace
(b) & \PROVED\ as cohomological correspondences.  In the Haran Hilbert state,
\(V_n=F_n^*\), and the self-dual metric normalization produces the
\(\Lambda(n)/\sqrt n\) coefficient. & No undecorated Cartier/Chow lift is
asserted.  The Hilbert adjoint in this model is not automatically the adjoint of
the odd nuclear quotient used in row (c). \\
\addlinespace
(c) & \PROVED\ as a cohomological/nuclear Lefschetz identity.  The even part
contains the two Tate directions and the odd part is the nuclear Fr\'echet
quotient \(H_-^0=H_-/ZH_{\cap}\). & It does not construct a positive Hilbert
completion of that quotient for which the nuclear involution becomes the Hilbert
adjoint and the nuclear character is preserved. \\
\bottomrule
\end{longtable}

The tempting implication
\begin{equation}
 V_n=F_n^*
 \quad\Longrightarrow\quad
 \pi_-(f^\vee)=\pi_-(f)^*
 \tag{A1}
\end{equation}
is therefore not among the proved results.  File v96 proves the algebraic scaling
adjoint
\[
 V_n^\dagger\phi_k=n^{-1}\phi_{k/n}
 \quad(n\mid k),
 \qquad U_n^*=n^{-1}U_{1/n}
\]
on its algebraic domain, and proves descent of the involution.  It leaves open the
positive Hilbert seminorm, closedness of \(ZH_{\cap}\), and preservation of the odd
nuclear character.  File v97 further proves that the existing positive
GNS/Gamma--Euler completions do not automatically descend to this quotient and
that a cyclic matrix coefficient is not the nuclear trace.

\medskip
\noindent\AUDIT\quad The even Tate superobject in row (c) is essential structural
data, but it is not by itself the missing positive metric on the odd quotient.
Using (A1) without constructing and checking that metric would insert the desired
positivity as an unproved assumption.

\section{New row (d): exact ledger of what survives the audit}

\subsection{Core exact identities}

\begin{enumerate}
\item[\PROVED] \textbf{G39 finite Gamma--Tate amplitude (v94).}
The exact finite-threshold calculation gives
\[
 \mathcal H_{\infty,0}^*\mathcal W_\infty=-B_Q
\]
and the combined divergence equals \(M_0h'\).  This repairs the earlier mismatch
between the moving harmonic residual and the differentiated metric.

\item[\PROVED] \textbf{Moving Tate decomposition and conservation (v117).}
The moving cell, its isometry, conservation identity, and full compression formula
are established.  The final physical innovation factorization is not established.

\item[\PROVED] \textbf{Cross-only G39 transport (v119).}
The physical cross incidence, arithmetic cross Gram, and word-column description
are exact.  The shorted inverse-metric domination is not a consequence of these
identities.

\item[\PROVED] \textbf{Positive Gamma source Gram (v120).}
The Gamma contribution has a positive Gram representation.  This does not imply
the required inverse-metric estimate, because compression and inversion do not
commute.

\item[\PROVED] \textbf{Causal Gamma state and leakage split (v121).}
Integer and noninteger labels separate orthogonally at the relevant level.  The
noninteger labels can form connected rational clusters; uniform unit domination
of those clusters is open.

\item[\PROVED] \textbf{Local Gamma--Tate cell bound (v122).}
A one-cell inequality with a positive local margin is established.  It cannot be
restarted independently at every cell: a connected cluster transports a
near-constant Gamma state across cells.

\item[\PROVED] \textbf{Exact correlated rational leakage (v123).}
The coefficient is
\[
 \ell_{N,k}(a,b)=
 \frac{1}{\sqrt{ab}(\log N)^k}
 \sum_{d\leq N/\max(a,b)}
 \frac{\Lambda(bd)\Lambda_k(ad)}{d}.
\]
This identifies data not determined by the diagonal Witt moments alone.
\end{enumerate}

\subsection{The uniform rational Gamma--Tate route}

File v124 formulates the exact unit-constant estimate
\begin{equation}
 \|L_Ne\|^2\leq \mathcal E_{\Gamma T,N}(e)
 \qquad\text{for every }N\text{ and every admissible }e.
 \tag{URGT_1}
\end{equation}
Its scalar expansion and its equivalence to the operator inequality in that model
are \PROVED.  The estimate itself is \OPEN.

The cascade
\begin{equation}
 (\mathrm{URGT}_1)
 \Longrightarrow
 \text{uniform shorted leakage domination}
 \Longrightarrow G40
 \Longrightarrow G
 \Longrightarrow \text{row (d)}
 \tag{C1}
\end{equation}
is a \CONDITIONAL\ sufficient route.  It is not proved that
\((\mathrm{URGT}_1)\) is necessary for G40.  Consequently, failure or incompletion
of this route would not kill row (d), and its subsidiary lemmas must not be called
the unique remaining steps of row (d).

Files v125--v130 analyze one compactness/minimizer strategy for
\((\mathrm{URGT}_1)\).  They provide fixed-\(N\) coercivity, all-depth local profile
information, and useful resolvent/compactness reductions.  They do not establish
uniform observability in the moving arithmetic holes where the weight is small.
That observability statement is \OPEN, but it is only the remaining gap of this
particular strategy, not the definition of G40 and not the sole logically possible
closure of row (d).

\section{Why the numerical endpoint is not the global closure}

The old row (d) proves genuine numerical and interval statements.  They are useful
tests of normalization, signs, and the first Schur cells.  Nevertheless,
\begin{equation}
 \text{one threshold or one finite compression is positive}
 \quad\not\Longrightarrow\quad
 P_T\mathcal A_TP_T\geq0\quad\text{for every }T.
 \tag{N1}
\end{equation}
At \(T=\tfrac12\log5\), separate positivity of a finite block and its complement
does not control their off-diagonal coupling.  That coupling is exactly what the
Schur complement measures.  Hence the ``infinitesimal decimal'' work is a valid
local certification and debugging tool, but it is not a proof of the uniform
theorem.

\section{Corrected terminal table}

\begin{longtable}{>{\raggedright\arraybackslash}p{0.31\textwidth}
                  >{\raggedright\arraybackslash}p{0.14\textwidth}
                  >{\raggedright\arraybackslash}p{0.45\textwidth}}
\toprule
Item & Status & Exact meaning \\
\midrule
\endfirsthead
\toprule
Item & Status & Exact meaning \\
\midrule
\endhead
Row (a) & \PROVED & Complete in its stated intrinsic
periodic--Deligne--nuclear category. \\
Row (b) & \PROVED & Complete as cohomological correspondences with its stated
Hilbert normalization; no undecorated Chow lift. \\
Row (c) & \PROVED & Complete as a nuclear/cohomological Lefschetz identity; no
positive Hilbert completion of the odd quotient is asserted. \\
Old row (d), \(0<T\leq\log2\) & \PROVED & Full Hilbert positivity on the certified
initial interval. \\
Old row (d), all \(T\) & \OPEN & The uniform Schur/capacity estimate is absent. \\
G39 & \PROVED & Exact finite Gamma--Tate amplitude/divergence identity. \\
Moving Tate conservation & \PROVED & Exact decomposition, isometry, and
compression identity. \\
G40 & \OPEN & The Hilbert positive factorization (F4), equivalently (F2)--(F3),
has not been proved for every threshold. \\
G41: \(J_T=I\) & \OPEN & Constructed and proved equivalent to G40; its value is
not proved. \\
G42: \(\Delta_T\geq0\) & \OPEN & Exact defect constructed; its sign is the
load-bearing inequality. \\
Positive Hilbert Rosati descent & \OPEN & Alternate closure route; the required
metric/quotient/trace compatibility is not constructed. \\
Correlated rational leakage domination & \OPEN & Direct missing unit domination
in the Gamma-state route. \\
\((\mathrm{URGT}_1)\) & \OPEN & Explicit sufficient uniform theorem, not a proved
fact and not shown necessary. \\
Uniform minimizer observability & \OPEN & Gap in one proposed proof of
\((\mathrm{URGT}_1)\), not a new definition of row (d). \\
Condition \(G\) & \OPEN & Equivalent positive statement not established. \\
New row (d), all \(T\) & \OPEN & Depends on a genuine proof of G40 or another
independent theorem implying the same positivity. \\
\bottomrule
\end{longtable}

\section{What is useful, and what should leave the critical path}

\subsection*{Retain in the main proof spine}
\begin{enumerate}
\item Rows (a), (b), and (c), with their category boundaries stated exactly.
\item The old row (d) definition, the certified interval \(T\leq\log2\), and the
      exact threshold Schur/capacity reduction.
\item The v94 G39 amplitude and the single frozen G40 equivalence
      (F1)--(F5).
\item The v117 and v119--v123 exact moving-Tate/Gamma/leakage identities, because
      they identify the physical source and the actual correlated obstruction.
\item v124 as a clearly labelled sufficient conjectural theorem if the authors
      choose to pursue the rational Gamma--Tate route.
\end{enumerate}

\subsection*{Move to an audit/no-go appendix}
The alternative machines in v95--v116 and the compactness experiments v125--v130
are valuable route audits, counterexample filters, and records of what cannot be
inferred.  They should not appear as successive ``almost final'' steps in the main
logical chain.  None presently supplies the missing sign.

\section{The three honest closure routes}

There are presently three noncircular targets.  Proving any one with all domain,
range, and convergence statements would be a real closure; merely defining it
would not.

\begin{enumerate}
\item \textbf{Direct G40.}  Prove (F2)--(F3), or equivalently construct the
Hilbert factorization (F4), from the explicit Gamma--Euler--Tate network.  This is
the shortest target and should remain the frozen definition of success.

\item \textbf{Uniform rational Gamma--Tate theorem.}  Prove
\((\mathrm{URGT}_1)\) and verify the already formulated conditional cascade (C1)
without changing constants or spaces.  The current minimizer-observability work is
one possible proof strategy, not part of the statement to be proved.

\item \textbf{Positive Rosati descent.}  Construct a positive Hilbert completion
of the odd Poisson quotient, prove the null subspace is closed, prove that
\(f\mapsto f^\vee\) becomes the Hilbert adjoint, and prove that the resulting
Hilbert trace equals the nuclear character of row (c).  Only after these checks can
the negative Hilbert--Schmidt square be used.  The existing files show that this
route carries load-bearing content rather than a formal transport.
\end{enumerate}

\section{Final audit conclusion}

The work has not been empty and the boundary has not merely been renamed.  G39,
the moving Tate state, exact conservation, the Gamma source Gram, the rational
cluster decomposition, and the correlated leakage coefficient are genuine new
mathematical constructions.  They explain precisely why simpler diagonal,
cellwise, or formal-adjoint closures fail.

But the current corpus contains no proof of the uniform positive inequality.
Accordingly, the only defensible final labels are
\[
 \boxed{
 \text{G39: \PROVED,}\qquad
 \text{G40/G41/G42: \OPEN,}\qquad
 G:\ \OPEN,\qquad
 \text{new row (d): \OPEN.}}
\]
This is the stable ledger.  Future work should be measured directly against (F3)
or (F4); no subsidiary reduction is to be advertised as the new unique terminal
gap unless its equivalence to G40 is independently proved.

\paragraph{Typographical audit.}
The v117 source contains a malformed \verb|\binom| token and a missing backslash
in one occurrence of \verb|\ker B_T|.  In v122, equation (7) appears to omit a
plus sign between the Gamma energy and the endpoint term; equation (6), its proof,
and subsequent use determine the intended sign.  These are editorial defects, not
mathematical closures or counterexamples.

\end{document}
